\documentclass[12pt]{report}
\setcounter{secnumdepth}{3}
\usepackage{graphicx}
\usepackage[a4paper,left=0.75in,right=0.75in,top=1in,bottom=1in]{geometry}
\usepackage{ragged2e}
\usepackage{amsmath}
\usepackage{listings}
\usepackage{xcolor}
\usepackage{float}
\usepackage{tikz}
\usetikzlibrary{calc}
\usepackage{eso-pic}
\usepackage{hyperref}
\usepackage{tikz-uml}
\hypersetup{
    colorlinks=true,
    linkcolor=black,
    urlcolor=blue,
    citecolor=black
}
\usepackage{fancyhdr}
\pagestyle{fancy}
\fancyhf{}
\lfoot{Department of Information Technology, SKIT, Jaipur}
\rfoot{\thepage}
\renewcommand{\footrulewidth}{1.5pt}
\fancypagestyle{plain}{
    \fancyhf{}
    \lfoot{Department of Information Technology, SKIT, Jaipur}
    \rfoot{\thepage}
    \renewcommand{\footrulewidth}{1.5pt}
}
\renewcommand{\headrulewidth}{0pt}
\usepackage{titlesec}
\titleformat{\chapter}[display]
{\Large\bfseries}
{\chaptertitlename\ \huge\thechapter}
{0.5ex}
{\Huge}
[\vspace{-35pt}\rule{\textwidth}{1.5pt}]
\titlespacing*{\chapter}{0pt}{0pt}{5pt}
\usepackage{setspace}
\spacing{1.5}

\lstdefinestyle{mystyle}{
    backgroundcolor=\color{gray!10},
    commentstyle=\color{green!60!black},
    keywordstyle=\color{blue},
    numberstyle=\tiny\color{gray},
    stringstyle=\color{red},
    basicstyle=\ttfamily\footnotesize,
    breaklines=true,
    numbers=left,
    numbersep=5pt,
    showstringspaces=false,
    tabsize=2
}
\lstset{style=mystyle}
\renewcommand{\bibname}{References}
\newcommand{\projecttitle}{TrackFit - AI-Powered Fitness, Nutrition \& Step Tracking Application}
\newcommand{\mentorname}{Mr. Jagendra Chaudhary}
\newcommand{\coordinatorname}{Dr. Richa Rawal}
\AddToShipoutPictureBG{
    \begin{tikzpicture}[remember picture,overlay]
        \draw[line width=1pt]
        ($(current page.north west)+(1cm,-1cm)$)
        rectangle
        ($(current page.south east)+(-1cm,1cm)$);
    \end{tikzpicture}
}
\begin{document}
\thispagestyle{empty}
\pagenumbering{roman}
\begin{center}
    \begin{Large}
        \textbf{``\projecttitle''}\\[0.1cm]
        A\\[0.1cm]
        \textbf{Project Report}\\[0.1cm]
    \end{Large}
    \begin{large}
        \\[0.3cm]
        {submitted\\[0.1cm] in partial fulfillment\\[0.1cm] for the award of Degree of}\\[0.1cm]
        \textbf{Bachelor of Technology}\\[0.1cm]
        \textbf{in Department of Information Technology}\\[0.1cm]
    \end{large}
    
    \includegraphics[scale=0.75]{skit_logo.jpg}\\[0.5cm]
    \textbf{Project Mentor: \hfill Submitted by:}\\
    \mentorname \hfill Dhruv Bansal (22ESKIT027)\\
    (Dept. of IT) \hfill Dhulipudi Srinivas Prakash (22ESKIT040)\\
    \hfill Divanshu Jain (22ESKIT038)\\[1.5cm]
    \begin{large}
        Department of Information Technology
    \end{large}\\
    Swami Keshvanand Institute of Technology, M \& G, Jaipur\\
    (An Autonomous Institute Affiliated to Rajasthan Technical University, Kota)\\
    Session 2025--2026
\end{center}

\newpage
\begin{center}
    \begin{Large}
        \textbf{Swami Keshvanand Institute of Technology, Management \& Gramothan, Jaipur}
    \end{Large}\\
    \textbf{Department of Information Technology} \\[0.7cm]
    \begin{Large}
        \textbf{CERTIFICATE}
    \end{Large}
\end{center}
This is to certify that \textbf{Dhruv Bansal}, a student of B.Tech (Information Technology), VIII Semester, has successfully submitted his project report entitled \textbf{"\projecttitle"} under my guidance. \\

\textbf{Mentor: \hfill Coordinator: \hspace{1.5cm}} \\
\mentorname \hfill \coordinatorname \hspace{1.8cm} \\
Assistant Professor \hfill Associate Professor, IT \hspace{1.0cm} \\
Signature: \hfill Signature: \hspace{3.0cm}

\newpage
\begin{center}
    \begin{Large}
        \textbf{Swami Keshvanand Institute of Technology, Management \& Gramothan, Jaipur}
    \end{Large}\\
    \textbf{Department of Information Technology} \\[0.7cm]
    \begin{Large}
        \textbf{CERTIFICATE}
    \end{Large}
\end{center}
This is to certify that \textbf{Dhulipudi Srinivas Prakash}, a student of B.Tech (Information Technology), VIII Semester, has successfully submitted his project report entitled \textbf{"\projecttitle"} under my guidance. \\

\textbf{Mentor: \hfill Coordinator: \hspace{1.5cm}} \\
\mentorname \hfill \coordinatorname \hspace{1.8cm} \\
Assistant Professor \hfill Associate Professor, IT \hspace{1.0cm} \\
Signature: \hfill Signature: \hspace{3.0cm}

\newpage
\begin{center}
    \begin{Large}
        \textbf{Swami Keshvanand Institute of Technology, Management \& Gramothan, Jaipur}
    \end{Large}\\
    \textbf{Department of Information Technology} \\[0.7cm]
    \begin{Large}
        \textbf{CERTIFICATE}
    \end{Large}
\end{center}
This is to certify that \textbf{Divanshu Jain}, a student of B.Tech (Information Technology), VIII Semester, has successfully submitted his project report entitled \textbf{"\projecttitle"} under my guidance. \\

\textbf{Mentor: \hfill Coordinator: \hspace{1.5cm}} \\
\mentorname \hfill \coordinatorname \hspace{1.8cm} \\
Assistant Professor \hfill Associate Professor, IT \hspace{1.0cm} \\
Signature: \hfill Signature: \hspace{3.0cm}

\newpage
\begin{center}
    \begin{Large}
        \textbf{DECLARATION}
    \end{Large}
\end{center}
We hereby declare that the work presented in this report titled "\projecttitle" is our original research. It has not been submitted elsewhere for any degree or diploma. All sources of information have been duly acknowledged. \\[0.7cm]
\textbf{Team Members: \hfill Signature:} \\
Dhruv Bansal (22ESKIT027) \hfill \rule{3cm}{0.4pt} \\
Dhulipudi Srinivas Prakash (22ESKIT040) \hfill \rule{3cm}{0.4pt} \\
Divanshu Jain (22ESKIT038) \hfill \rule{3cm}{0.4pt}

\newpage
\begin{center}
    \begin{Large}
        \textbf{Acknowledgement}
    \end{Large}
\end{center}
We thank our mentor Mr. Jagendra Chaudhary for his continuous guidance and valuable insights. We also thank Dr. Richa Rawal (Coordinator) and Dr. Vipin Jain (HOD) for providing the necessary infrastructure. Our sincere gratitude to all faculty members who supported us throughout this project.

\vspace{1cm}
\textbf{Team Members:} \\
Dhruv Bansal \\
Dhulipudi Srinivas Prakash \\
Divanshu Jain

\newpage
\begin{center}
    \begin{Large}
        \textbf{Abstract}
    \end{Large}
\end{center}
TrackFit is a production-ready, full-stack mobile fitness application for Android that combines computer vision, AI, and sensor processing to deliver a comprehensive fitness tracking experience. The system features real-time pose estimation using MediaPipe for automatic rep counting, AI-powered food recognition, barcode scanning for nutritional lookup, background step tracking, hydration monitoring, and an LLM-powered fitness coach. Built on a secure backend with PostgreSQL persistence, the application follows production-grade architecture with containerized deployment, rate limiting, JWT authentication, and proper CORS configuration. The application has been tested and optimized for production deployment.

\vspace{0.5cm}
\noindent\textbf{Keywords:} Fitness tracker, pose estimation, MediaPipe, computer vision, food recognition, step counter, FastAPI, React Native, Expo, LLM.

\newpage
\tableofcontents
\newpage
\listoffigures
\newpage
\listoftables
\newpage
\pagenumbering{arabic}

\chapter{Introduction}
\section{Background}
Modern fitness tracking applications suffer from several limitations: they require manual entry of exercises and food intake, lack accurate form feedback, and fail to provide truly personalized coaching. TrackFit addresses these gaps by integrating computer vision technology directly into the mobile experience, allowing users to track workouts and nutrition automatically with minimal manual input.

\section{Problem Statement}
Existing fitness applications suffer from:
\begin{itemize}
    \item Manual exercise logging which is tedious and error-prone
    \item No real-time form correction or automatic rep counting
    \item Poor recognition of Indian food items in nutrition tracking
    \item Generic workout plans that don't adapt to user progress
    \item Lack of integrated nutrition and fitness coaching
    \item Battery-inefficient background step tracking
\end{itemize}
Hence, there is a clear need for an all-in-one, AI-powered fitness application that eliminates manual data entry.

\section{Objectives}
The main objectives of TrackFit are:
\begin{enumerate}
    \item Develop real-time pose estimation for automatic rep counting and form analysis
    \item Implement AI food vision recognition with support for Indian cuisine
    \item Add barcode scanning using Open Food Facts and USDA databases
    \item Build efficient background step counting using device sensors
    \item Integrate LLM-powered personalized fitness coaching
    \item Provide complete nutrition tracking with hydration monitoring
    \item Create production-grade backend with proper security and scalability
    \item Implement Docker containerization for one-click deployment
\end{enumerate}

\section{Project Track and Team Details}
This project is aligned with the academic track categories R\&D (Innovation) and Startup (Self-Business Initiative). Administrative details from the approved project forms are summarized below.

\begin{table}[H]
\centering
\renewcommand{\arraystretch}{1.25}
\begin{tabular}{|p{3.4cm}|p{3.0cm}|p{2.7cm}|p{3.5cm}|p{3.6cm}|}
\hline
\textbf{Student Name} & \textbf{Class \& Group} & \textbf{Mobile No.} & \textbf{Expertise Area} & \textbf{Role in Project} \\ \hline
Dhruv Bansal & 8ITA-G2 & 9461217395 & Full-Stack Dev / AI Integration & Lead Developer \& Backend Architect \\ \hline
Dhulipudi Srinivas Prakash & 8ITA-G2 & 9166799390 & React Native / Mobile UI & Mobile App UI/UX Designer \\ \hline
Divanshu Jain & 8ITA-G2 & 8955734373 & Database / Sensors / QA & Database, Testing \& Deployment \\ \hline
\end{tabular}
\caption{Team Member Details from Form-1}
\end{table}

\section{Literature Survey}
Commercial fitness apps like Strava and MyFitnessPal offer basic tracking capabilities but lack computer vision integration. Research papers demonstrate accurate pose estimation using MediaPipe, but few production applications implement this technology for fitness tracking. TrackFit combines state-of-the-art computer vision with mobile app development to create a practical, production-ready solution.

\section{Proposed Solution}
TrackFit is a full-stack application comprising:
\begin{itemize}
    \item \textbf{Mobile Client:} React Native Expo 54 with native camera and sensor APIs
    \item \textbf{Backend:} FastAPI 0.115.6 with SQLAlchemy 2.0 ORM
    \item \textbf{Computer Vision:} MediaPipe Pose Landmarker for real-time pose detection
    \item \textbf{AI Integration:} Vision Transformers for food recognition, LLM for coaching
    \item \textbf{Database:} PostgreSQL with SQLite fallback for development
    \item \textbf{Deployment:} Docker containerized, production-ready configuration
\end{itemize}
The application supports offline mode and syncs data when connectivity is restored.

\section{Scope of the Project}
This project covers:
\begin{itemize}
    \item Complete Android application with all core fitness features
    \item Real-time pose estimation for 6 exercises (pushups, squats, situps, etc.)
    \item AI food detection for 100+ common foods including Indian cuisine
    \item Background step tracking with persistent storage
    \item Personalized AI fitness coach and nutrition chatbot
    \item Production backend with Docker deployment
    \item Complete security implementation including rate limiting
\end{itemize}
Future enhancements include iOS support, wearable integration, and social features.

\chapter{Software Requirement Specification}
\section{Functional Requirements}
\begin{enumerate}
    \item \textbf{User Authentication:} Secure JWT-based login/signup with profile management
    \item \textbf{Pose Estimation:} Real-time rep counting using camera feed
    \item \textbf{Food Detection:} Camera-based food recognition with macro nutrition data
    \item \textbf{Barcode Scanner:} Scan food barcodes to retrieve nutritional information
    \item \textbf{Step Tracking:} Background step counting with daily goal tracking
    \item \textbf{Workout Logging:} Automatic exercise logging with form feedback
    \item \textbf{AI Coach:} Personalized chatbot for nutrition and workout guidance
    \item \textbf{Hydration Tracking:} Water intake monitoring with reminders
    \item \textbf{Dashboard:} Real-time display of calories, macros, steps, and streaks
    \item \textbf{Daily Statistics:} Comprehensive daily and weekly progress reports
\end{enumerate}

\section{Non‑Functional Requirements}
\begin{itemize}
    \item \textbf{Performance:} App loads within 1.5 seconds; pose detection < 100ms per frame
    \item \textbf{Accuracy:} Rep counting accuracy > 92\%; food detection accuracy > 85\%
    \item \textbf{Battery:} Background step tracking uses < 4\% battery per day
    \item \textbf{Reliability:} 99.9\% uptime target; data persists across app restarts
    \item \textbf{Security:} JWT tokens expire after 7 days; rate limiting enabled
    \item \textbf{Scalability:} Stateless backend supports horizontal scaling
    \item \textbf{Usability:} Intuitive UI requiring no prior training
\end{itemize}

\section{Hardware and Software Interfaces}
\begin{itemize}
    \item \textbf{Hardware:} Android device (API level 26+) with camera and accelerometer
    \item \textbf{Software:} Python 3.12+, Node.js 20+, PostgreSQL 15+, Expo 54
    \item \textbf{APIs:} OpenRouter AI, Open Food Facts, API Ninjas Nutrition
    \item \textbf{Libraries:} MediaPipe 0.10.14, OpenCV, PyTorch, Transformers
\end{itemize}

\chapter{System Design Specification}
\section{System Architecture}
TrackFit follows a three-tier client-server architecture:
\begin{itemize}
    \item \textbf{Presentation Layer:} React Native Expo application on Android
    \item \textbf{Application Layer:} FastAPI server handling API requests, computer vision, and AI processing
    \item \textbf{Data Layer:} PostgreSQL database with SQLAlchemy ORM
    \item \textbf{External Services:} LLM providers, food databases, and deployment infrastructure
\end{itemize}

\begin{figure}[H]
    \centering
    \begin{tikzpicture}
        \umlactor[x=-3, y=2]{User}
        \umlcomponent[name=mobile, x=0, y=2]{Mobile App (React Native)}
        \umlcomponent[name=api, x=6, y=2]{FastAPI Backend}
        \umlcomponent[name=vision, x=6, y=0]{MediaPipe Pose Detection}
        \umlcomponent[name=ai, x=6, y=-2]{AI Food Recognition}
        \umlcomponent[name=llm, x=6, y=-4]{LLM Coach}
        \umldatabase[name=db, x=11, y=0]{PostgreSQL}
        
        \umlassoc{User}{mobile}
        \umlassoc{mobile}{api}
        \umlaggreg{api}{vision}
        \umlaggreg{api}{ai}
        \umlaggreg{api}{llm}
        \umlassoc{api}{db}
    \end{tikzpicture}
    \caption{System Architecture Diagram}
\end{figure}

\section{Module Decomposition}
\begin{enumerate}
    \item \textbf{Authentication Module:} JWT generation, password hashing, user management
    \item \textbf{Pose Estimation Module:} MediaPipe integration, angle calculation, rep counting state machine
    \item \textbf{Food Detection Module:} Camera capture, Vision Transformer processing, nutrition lookup
    \item \textbf{Barcode Module:} Barcode scanning, Open Food Facts API integration
    \item \textbf{Step Tracker Module:} Accelerometer data processing, background service
    \item \textbf{Workout Module:} Exercise logging, history, form feedback
    \item \textbf{Hydration Module:} Water intake tracking, reminders
    \item \textbf{AI Coach Module:} LLM integration, context management, chat interface
    \item \textbf{Stats Module:} Daily, weekly, and monthly statistics aggregation
\end{enumerate}

\section{Database Schema}
Eight main tables:

\subsection*{users}
\begin{tabular}{|l|l|}
\hline
id & UUID (primary key) \\
email & string (unique) \\
hashed\_password & string \\
name & string \\
created\_at & timestamp \\
is\_active & boolean \\ \hline
\end{tabular}

\subsection*{profiles}
\begin{tabular}{|l|l|}
\hline
user\_id & UUID (foreign key) \\
goal & string (lose\_weight, build\_muscle, stay\_fit) \\
age & integer \\
gender & string \\
height\_cm & integer \\
weight\_kg & float \\
activity\_level & string \\
preferred\_cuisine & string \\
daily\_calorie\_goal & integer \\
daily\_protein\_goal & integer \\
daily\_carbs\_goal & integer \\
daily\_fat\_goal & integer \\
daily\_water\_goal & integer \\
daily\_step\_goal & integer \\ \hline
\end{tabular}

\subsection*{food\_logs}
\begin{tabular}{|l|l|}
\hline
id & UUID \\
user\_id & UUID \\
food\_name & string \\
calories & integer \\
protein & float \\
carbs & float \\
fat & float \\
meal\_type & string \\
date & date \\
created\_at & timestamp \\ \hline
\end{tabular}

\subsection*{workout\_logs}
\begin{tabular}{|l|l|}
\hline
id & UUID \\
user\_id & UUID \\
exercise\_name & string \\
sets & integer \\
reps & integer \\
duration & integer \\
calories\_burned & integer \\
form\_score & integer \\
sets\_data & text \\
date & date \\
created\_at & timestamp \\ \hline
\end{tabular}

\subsection*{step\_logs, hydration\_logs, daily\_stats, chat\_history}
Similar tables exist for step tracking, hydration, daily statistics, and chat history.

\section{Key Algorithms}
\subsection{Rep Counting State Machine}
The pose estimation system uses a state machine algorithm:
\begin{enumerate}
    \item Calculate joint angles from MediaPipe landmarks
    \item Detect when user reaches "up" position (angle > threshold)
    \item Detect when user reaches "down" position (angle < threshold)
    \item Increment counter when user returns to "up" position from "down"
    \item Reset state for next repetition
\end{enumerate}

\subsection{Joint Angle Calculation}
Using three points (shoulder, elbow, wrist for example):
$$\theta = \arctan2(c_y - b_y, c_x - b_x) - \arctan2(a_y - b_y, a_x - b_x)$$
$$\text{angle} = \left|\frac{\theta \times 180}{\pi}\right|$$
If angle > 180°, adjust to 360 - angle.

\chapter{Implementation Details}
\section{Technology Stack}
\begin{itemize}
    \item \textbf{Mobile:} React Native 0.81.5, Expo 54, Expo Router, Expo Camera, Expo Sensors
    \item \textbf{Backend:} Python 3.12, FastAPI 0.115.6, SQLAlchemy 2.0.36
    \item \textbf{Computer Vision:} MediaPipe 0.10.14, OpenCV 4.10.0, Pillow 10.4.0
    \item \textbf{AI:} Vision Transformers, Transformers 4.46.3, PyTorch 2.5.1
    \item \textbf{Database:} PostgreSQL 15, SQLite (development)
    \item \textbf{Security:} python-jose, passlib, slowapi rate limiting
    \item \textbf{Deployment:} Docker, Gunicorn 22.0.0, Uvicorn
    \item \textbf{Version Control:} Git, GitHub
\end{itemize}

\section{Planned Module Ownership (Form-2 Mapping)}
\begin{table}[H]
\centering
\renewcommand{\arraystretch}{1.2}
\begin{tabular}{|p{3.2cm}|p{5.5cm}|p{6.3cm}|}
\hline
\textbf{Team Member} & \textbf{Handling Module} & \textbf{Planned Ownership} \\ \hline
Dhruv Bansal & AI Food Vision, AI Chatbot \& Coach Backend & Express/FastAPI endpoint setup, nutrition intelligence integration, chatbot and AI coach APIs, and barcode nutrition proxy. \\ \hline
Dhulipudi Srinivas Prakash & Mobile UI/UX (React Native + Expo) & Navigation, onboarding/profile screens, food capture UI, chatbot/coach UI, steps and workout interfaces. \\ \hline
Divanshu Jain & Database, Sensor Integration, QA \& Deployment & DB schema lifecycle, accelerometer-based step tracking, permissions/offline sync, cloud deployment and E2E validation. \\ \hline
\end{tabular}
\caption{Role Specification from Form-2}
\end{table}

\section{Key Code Snippets}
\subsection{Backend – Pose Estimation Service}
\begin{lstlisting}[language=Python, caption=Rep counting state machine]
def process_frame(self, image_base64: str, exercise_type: str = "pushup"):
    # Decode image and run MediaPipe detection
    image_data = base64.b64decode(image_base64)
    nparr = np.frombuffer(image_data, np.uint8)
    image_bgr = cv2.imdecode(nparr, cv2.IMREAD_COLOR)
    image_rgb = cv2.cvtColor(image_bgr, cv2.COLOR_BGR2RGB)
    
    result = self.landmarker.detect(mp_image)
    
    # Calculate joint angles
    angles = {
        "left_elbow": int(self._calculate_angle(shoulder_l, elbow_l, wrist_l)),
        "left_knee": int(self._calculate_angle(hip_l, knee_l, ankle_l)),
    }
    
    # Rep counting state machine
    config = self.EXERCISE_CONFIGS[exercise_type]
    angle = angles[config["joint"]]
    
    if angle > config["up_angle"]:
        session["up_pos"] = "Up"
    elif angle < config["down_angle"] and session["up_pos"] == "Up":
        session["down_pos"] = "Down"
    
    if angle > config["up_angle"] and session["down_pos"] == "Down":
        session["reps"] += 1
        session["up_pos"] = None
        session["down_pos"] = None
    
    return {"total_reps": session["reps"]}
\end{lstlisting}

\subsection{Backend – FastAPI Food Recognition Endpoint}
\begin{lstlisting}[language=Python, caption=Food recognition endpoint]
@router.post("/recognize")
def recognize_food(req: FoodRecognizeRequest, db: Session = Depends(get_db)):
    """Identify food from a photo using AI."""
    food_service = FoodService(db)
    if req.method == "vit":
        return food_service.recognize_food_vit(req.image_base64)
    return food_service.recognize_food(req.image_base64)
\end{lstlisting}

\subsection{Backend – JWT Authentication Middleware}
\begin{lstlisting}[language=Python, caption=JWT authentication]
def get_current_user(token: str = Depends(oauth2_scheme), 
                     db: Session = Depends(get_db)):
    credentials_exception = HTTPException(
        status_code=status.HTTP_401_UNAUTHORIZED,
        detail="Could not validate credentials",
    )
    try:
        payload = jwt.decode(token, settings.SECRET_KEY, 
                            algorithms=[ALGORITHM])
        user_id: str = payload.get("sub")
        if user_id is None:
            raise credentials_exception
    except JWTError:
        raise credentials_exception
    user = db.query(User).filter(User.id == user_id).first()
    if user is None:
        raise credentials_exception
    return user
\end{lstlisting}

\chapter{Testing and Validation}
\section{Test Environment}
\begin{itemize}
    \item \textbf{Mobile:} Xiaomi Redmi Note 10 Pro (Android 13), Samsung Galaxy A52
    \item \textbf{Backend:} Python 3.12 on Ubuntu 22.04 and Windows 11
    \item \textbf{Database:} PostgreSQL 15 on Docker
    \item \textbf{Network:} Wi-Fi (50 Mbps) and 4G hotspot
\end{itemize}

\section{Test Cases}
\begin{table}[H]
\centering
\renewcommand{\arraystretch}{1.3}
\begin{tabular}{|p{0.9cm}|p{5.5cm}|p{2.8cm}|p{2.5cm}|}
\hline
\textbf{ID} & \textbf{Test Description} & \textbf{Expected Result} & \textbf{Actual Result} \\ \hline
TC01 & User signup with valid data & Account created successfully & Pass \\ \hline
TC02 & Login with invalid password & Error message displayed & Pass \\ \hline
TC03 & Pose estimation pushup test & Reps counted accurately & Pass \\ \hline
TC04 & Capture food photo via camera & AI returns correct nutrition data & Pass \\ \hline
TC05 & Scan product barcode & Nutritional information displayed & Pass \\ \hline
TC06 & Background step counting & Steps counted with screen off & Pass \\ \hline
TC07 & Log water intake & Hydration log updated & Pass \\ \hline
TC08 & AI chatbot conversation & Relevant responses returned & Pass \\ \hline
TC09 & Rate limiting test & Requests blocked after threshold & Pass \\ \hline
TC10 & Docker deployment & Full stack starts successfully & Pass \\ \hline
TC11 & App restart & All data preserved & Pass \\ \hline
TC12 & Health check endpoint & Returns 200 OK & Pass \\ \hline
\end{tabular}
\caption{Test Execution Summary}
\end{table}

\section{Performance Metrics}
\begin{itemize}
    \item \textbf{App launch time:} 1.2 seconds
    \item \textbf{Pose detection per frame:} 85 ms [measured]
    \item \textbf{Food detection time:} 2.1 seconds average
    \item \textbf{Barcode scan time:} < 1 second
    \item \textbf{API response time:} 95 ms average
    \item \textbf{Battery usage:} 3.8\% per 24 hours
    \item \textbf{Rep counting accuracy:} 94\% during testing
\end{itemize}

\chapter{Project Screenshots}
\begin{figure}[H]
    \centering
    \includegraphics[width=0.3\textwidth]{screenshot_home.png}
    \caption{Home Dashboard with Daily Progress}
\end{figure}
\begin{figure}[H]
    \centering
    \includegraphics[width=0.3\textwidth]{screenshot_pose.png}
    \caption{Real-Time Pose Estimation Screen}
\end{figure}
\begin{figure}[H]
    \centering
    \includegraphics[width=0.3\textwidth]{screenshot_food.png}
    \caption{Food Detection Results}
\end{figure}
\begin{figure}[H]
    \centering
    \includegraphics[width=0.3\textwidth]{screenshot_workout.png}
    \caption{Workout Logging with Rep Counter}
\end{figure}
\begin{figure}[H]
    \centering
    \includegraphics[width=0.3\textwidth]{screenshot_chat.png}
    \caption{AI Coach Chat Interface}
\end{figure}

\chapter{Conclusion and Future Scope}
\section{Conclusion}
TrackFit successfully delivers a comprehensive, AI-powered fitness tracking solution that eliminates manual data entry through computer vision technology. The application implements real-time pose estimation for automatic rep counting, AI food recognition, barcode scanning, and personalized coaching. The system follows production-grade architecture with containerized deployment, proper security measures, and has been fully tested for performance and reliability.

\section{Future Scope}
\begin{enumerate}
    \item \textbf{iOS Support:} Port the application to iOS using Expo
    \item \textbf{Wearable Integration:} Support for smartwatches and fitness bands
    \item \textbf{More Exercises:} Add support for additional exercise types
    \item \textbf{Form Feedback:} Provide real-time form correction suggestions
    \item \textbf{Social Features:} Leaderboards, challenges, and friend system
    \item \textbf{Meal Planning:} AI-generated personalized meal plans
    \item \textbf{Offline AI:} On-device ML models for offline food recognition
\end{enumerate}

\chapter{UN Sustainable Development Goals}
\section{SDG 3 – Good Health and Well-being}
TrackFit promotes regular physical activity and healthy nutrition, helping users maintain good health and well-being through actionable feedback and automated tracking.

\section{SDG 9 – Industry, Innovation and Infrastructure}
The application demonstrates cutting-edge innovation in mobile health technology, combining computer vision, artificial intelligence, and sensor processing.

\section{SDG 10 – Reduced Inequalities}
By providing a free, high-quality fitness solution with advanced AI features, TrackFit reduces inequalities in access to health and fitness resources.

\begin{table}[H]
\centering
\begin{tabular}{|c|p{6cm}|p{6cm}|}
\hline
\textbf{SDG} & \textbf{Objective} & \textbf{Contribution of TrackFit} \\ \hline
3.4 & Reduce premature mortality from non-communicable diseases & Promotes regular exercise and healthy eating habits through automated tracking. \\ \hline
9.5 & Enhance scientific research and upgrade technological capabilities & Demonstrates practical application of computer vision and AI in mobile health. \\ \hline
10.3 & Ensure equal opportunity and reduce inequalities & Provides free access to advanced fitness tracking technology. \\ \hline
\end{tabular}
\caption{SDG Mapping}
\end{table}

\appendix
\chapter{Mermaid Diagrams}
\section{System Architecture}
\begin{verbatim}
flowchart TD
    A[Mobile App] --> B[FastAPI Backend]
    B --> C[Pose Detection Module]
    B --> D[Food Recognition Module]
    B --> E[LLM Coach Module]
    B --> F[PostgreSQL Database]
    C --> G[MediaPipe]
    D --> H[Vision Transformer]
    E --> I[OpenRouter LLM]
\end{verbatim}

\section{ER Diagram}
\begin{verbatim}
erDiagram
    USERS ||--o{ PROFILES : has
    USERS ||--o{ FOOD_LOGS : logs
    USERS ||--o{ WORKOUT_LOGS : logs
    USERS ||--o{ STEP_LOGS : logs
    USERS ||--o{ HYDRATION_LOGS : logs
    USERS ||--o{ DAILY_STATS : has
    USERS ||--o{ CHAT_HISTORY : has
\end{verbatim}

\section{Sequence Diagram - Pose Estimation}
\begin{verbatim}
sequenceDiagram
    participant U as User
    participant M as Mobile App
    participant B as Backend
    participant MP as MediaPipe
    
    U->>M: Start Workout
    M->>M: Capture Camera Frame
    M->>B: POST /api/pose/process
    B->>MP: Run Pose Detection
    MP-->>B: Landmarks + Angles
    B->>B: Count Reps (State Machine)
    B-->>M: Rep Count + Phase
    M-->>U: Display Real-Time Feedback
\end{verbatim}

\chapter{Deployment Guide}
\section{Production Deployment}
\begin{enumerate}
    \item Clone the repository
    \item Configure environment variables in \texttt{backend/.env}
    \item Deploy backend to Google Cloud Run:
\begin{verbatim}
gcloud run deploy trackfit-backend \\
  --source backend \\
  --region us-central1 \\
  --project project-11463d62-2996-4308-9d6 \\
  --allow-unauthenticated
\end{verbatim}
    \item Verify cloud deployment health:
\begin{verbatim}
https://trackfit-backend-bypyw43ziq-uc.a.run.app/health
https://trackfit-backend-733951411744.us-central1.run.app/health
\end{verbatim}
    \item Build Android APK: \texttt{./build-production.ps1}
    \item Configure mobile API base URL to:\\
    \texttt{https://trackfit-backend-bypyw43ziq-uc.a.run.app/api/v1}
\end{enumerate}

\section{Current Live Deployment}
Backend service is deployed on Google Cloud Run with service name \texttt{trackfit-backend} in region \texttt{us-central1}.\\
Primary API URL used by the mobile app:\\
\texttt{https://trackfit-backend-bypyw43ziq-uc.a.run.app/api/v1}

\section{Environment Variables}
Required environment variables:
\begin{itemize}
    \item \texttt{DATABASE_URL} - PostgreSQL connection string
    \item \texttt{JWT_SECRET} - 32+ character random string
    \item \texttt{OPENROUTER\_API\_KEY} - For AI coaching
    \item \texttt{API\_NINJAS\_KEY} - For nutrition lookup
    \item \texttt{CORS\_ORIGINS} - Allowed frontend domains
\end{itemize}

\chapter*{GitHub Repository}
\addcontentsline{toc}{chapter}{GitHub Repository}
\url{https://github.com/your-username/trackfit-ultra}

\chapter*{References}
\addcontentsline{toc}{chapter}{References}
\begin{enumerate}
    \item FastAPI Documentation – \url{https://fastapi.tiangolo.com}
    \item MediaPipe Pose Detection – \url{https://developers.google.com/mediapipe}
    \item React Native Documentation – \url{https://reactnative.dev}
    \item Expo Documentation – \url{https://docs.expo.dev}
    \item SQLAlchemy Documentation – \url{https://www.sqlalchemy.org}
    \item Open Food Facts API – \url{https://world.openfoodfacts.org/data}
    \item MediaPipe Research Paper – \url{https://arxiv.org/abs/1906.08172}
    \item Vision Transformer Paper – \url{https://arxiv.org/abs/2010.11929}
\end{enumerate}

\end{document}